% !Mode:: "TeX:UTF-8"

\begin{publication}
\noindent\textbf{（一）发表的相关论文}

\begin{publist}
\item 一作. AER-Bench: A Benchmark for Emotion Perception and Reasoning in Fine Art[J]. IEEE Transactions on Affective Computing (\textbf{TAC}), 2026, pp.1-13. （EI, SCI, DOI: 10.1109/TAFFC.2026.3730731, IF: 11.3, 中科院1区TOP, JCR1区, CCF-B类期刊, 第5章）
\item Sicheng Zhao, \textbf{Xi Chen(共一)}, Hongxun Yao, Haosen Yang, Yanhao Zhang, Sheng Jin, Xiatian Zhu, Haonan Lu, Kui Jiang, Guiguang Ding. CMPF: Harmonizing Cross-Model Prior Fusion for Open-Vocabulary Segmentation[J]. International Journal of Computer Vision (\textbf{IJCV}), 2026, 134. （EI, SCI: 001779020700002, DOI: 10.1007/s11263-026-02886-0, IF: 10.3, 中科院2区，JCR1区, CCF-A类期刊, 第4章）
\item \textbf{Xi Chen}, Hongxun Yao, Sicheng Zhao, Jiankun Zhu, Jing Jiang, Kui Jiang. SfMamba: Efficient Source-Free Domain Adaptation via Selective Scan Modeling[J]. Expert Systems with Applications (\textbf{ESWA}), 2026, 324, pp.132421. （EI, SCI: 001759382000001, DOI: 10.1016/j.eswa.2026.132421, IF: 9.4, 中科院1区TOP, JCR1区, CCF-B类期刊，第2章）
\item \textbf{Xi Chen}, Haosen Yang, Huicong Zhang, Hongxun Yao and Xiatian Zhu. Uncertainty-aware Pseudo-label Filtering for Source-Free Unsupervised Domain Adaptation[J]. Neurocomputing, 2024, 575, pp.127190.（EI, SCI: 001172429800001, DOI: 10.1016/j.neucom.2023.127190, IF: 6.7, 中科院2区, JCR1区, CCF-C类期刊, 第2章）
\item \textbf{Xi Chen}, Hongxun Yao, Zhaopan Xu, Kui Jiang. DAM: Dual Active Learning with Multimodal Foundation Model for Source-Free Domain Adaptation[C]. IEEE International Conference on Acoustics, Speech and Signal Processing (\textbf{ICASSP}), 2026, pp. 5841-5845. （EI, DOI: 10.1109/ICASSP55912.2026.11463823, CCF-B类会议，第3章）
\end{publist}


\noindent\textbf{（二）在投的相关论文}

Xi Chen, Hongxun Yao, Zhaopan Xu, Sheng Jin, Sicheng Zhao, Kui Jiang. Data-Active Multimodal Fusion for Source-Free Domain Adaptation[J]. Information Fusion, 2026. (EI, SCI, IF: 17.4, 中科院1区TOP, JCR1区). （第3章）

\noindent\textbf{（三）已获得的专利}

% ===== 以下为非盲审原始内容，保留备用 =====
 姚鸿勋，陈希，张慧琮. 一种基于多基础大模型的开放词表分割方法. 发明专利, 2025.

\noindent\textbf{（四）参与的科研项目及获奖情况}

姚鸿勋等。农场环境自适应感知与认知技术研究，科技创新2030“新一代人工智能”重大项目。课题编号：2021ZD0110901。

% ===== 以下为非盲审原始内容，保留备用 =====
% \item 姚鸿勋等。农场环境自适应感知与认知技术研究，科技创新2030“新一代人工智能”重大项目。课题编号：2021ZD0110901。
%\vfill
%\hangafter=1\hangindent=2em\noindent
%\setlength{\parindent}{2em}
\end{publication}
